%====================================================================
\section{Giải mã và mã hóa file PLY}\label{s:codec}
%====================================================================
\src{io/codec.py: SplatCodec.decode / _decode_block / encode}
Một gaussian gồm tâm $\mu$, màu $c$, độ đục $\alpha$, kích thước $s$, quay $q$. File lưu giá trị ``thô'':

\begin{center}\small
\begin{tabular}{@{}lll@{}}\toprule
Thuộc tính & Giải mã & Lý do\\\midrule
\code{x,y,z} & $\mu=(x,y,z)$ & \\
\code{f_dc_0..2} & $c = 0.5 + C_0\,f_{dc}$,\; $C_0=\tfrac{1}{2\sqrt\pi}=0.28209479$ & $C_0=Y_{00}$ (SH bậc 0)\\
\code{opacity} & $\alpha = 1/(1+e^{-\mathrm{opacity}})$ & tối ưu logit để $\alpha\in(0,1)$\\
\code{scale_i} & $s_i = \exp(\mathrm{scale}_i)$ & tối ưu log để $s>0$\\
\code{rot_0..3} & $q\leftarrow q/\|q\|$; nếu $\|q\|=0$ thì $(1,0,0,0)$ & quaternion phải đơn vị\\\bottomrule
\end{tabular}
\end{center}

\textbf{Mã hóa ngược:} $f_{dc}=(c-0.5)/C_0$, $\mathrm{scale}=\ln\max(s,10^{-12})$, $\mathrm{opacity}=\ln\frac{a}{1-a}$ với $a=\mathrm{clip}(\alpha,10^{-6},1-10^{-6})$ để tránh $\ln 0$.

\textbf{Point cloud} (không có \code{scale_0}): bán kính mỗi điểm $=\tfrac12 (V/N)^{1/3}$, với $V$ là tích ba cạnh hộp bao. Đó là ước lượng khoảng cách trung bình giữa các điểm. Màu \code{uchar} chia 255.

\textbf{Dung lượng:} $3+3+3+45+1+3+4=62$ float32 $=248$ byte/gaussian (SH bậc 3: 15 hệ số $\times$ 3 kênh $=45$).

%====================================================================
\section{Biến đổi layer}\label{s:layer}
%====================================================================
\src{core/transform.py, core/layer.py, core/math3d.py}

\subsection{Ma trận quay từ góc Euler (\code{euler_to_matrix})}
\[
R_x=\begin{bmatrix}1&0&0\\0&c_x&-s_x\\0&s_x&c_x\end{bmatrix},\;
R_y=\begin{bmatrix}c_y&0&s_y\\0&1&0\\-s_y&0&c_y\end{bmatrix},\;
R_z=\begin{bmatrix}c_z&-s_z&0\\s_z&c_z&0\\0&0&1\end{bmatrix},\;
R=R_zR_yR_x
\]
Với vector cột, $X$ được áp dụng trước, rồi $Y$, rồi $Z$.

\subsection{Ngược lại (\code{matrix_to_euler})}
Từ $R=R_zR_yR_x$ có $R_{20}=-\sin y$, nên $y=\operatorname{asin}(-R_{20})$ (kẹp $[-1,1]$ chống lỗi làm tròn), $x=\operatorname{atan2}(R_{21},R_{22})$, $z=\operatorname{atan2}(R_{10},R_{00})$. Khi $|R_{20}|\ge0.999999$ ($y=\pm90^\circ$, \textbf{khóa gimbal}) thì $x,z$ không xác định riêng, đặt $x=0$, $z=\operatorname{atan2}(-R_{01},R_{11})$.

\subsection{Công thức Rodrigues (\code{axis_angle_matrix})}
Quay góc $a$ quanh trục đơn vị $(x,y,z)$, $c=\cos a$, $s=\sin a$, $t=1-c$:
\[
R=\begin{bmatrix}tx^2+c & txy-sz & txz+sy\\ txy+sz & ty^2+c & tyz-sx\\ txz-sy & tyz+sx & tz^2+c\end{bmatrix}
= I + s[k]_\times + t[k]_\times^2 .
\]

\subsection{Ma trận thế giới của layer (\code{Layer.world_transform / world_matrix})}
\[
p_w = s\,A\,p_f + t,\quad A=R(\mathrm{euler})\,F,\quad F=\mathrm{diag}(-1,-1,1),\quad
\Sig_w = s^2 A\Sig_f A^{\top},\quad
M=\begin{bmatrix}sA & t\\ 0 & 1\end{bmatrix}.
\]
\begin{itemize}
\item $F$ là phép \textbf{xoay $180^\circ$ quanh $Z$} ($\det=+1$, không phải phản xạ). Nó đưa dữ liệu COLMAP (Y hướng xuống) về Y hướng lên.
\item Chỉ $M$ thay đổi (dữ liệu gaussian giữ nguyên), nên kéo/xoay/zoom layer là tức thời. $M$ được gửi lên GPU thành 4 texel (\code{splat_renderer.set_transforms}, chuyển vị để ra 4 cột).
\end{itemize}

\subsection{Xoay quanh điểm pivot (\code{Transform.rotate_about(Δ, c)})}
$R'=\Delta R$, $\;t'=\Delta(t-c)+c$ (dời pivot về gốc, xoay, dời lại). Góc Euler mới lấy bằng \code{matrix_to_euler}$(\Delta R)$.

\subsection{Xuất file (\code{Layer._export_block})}
Đổi từ thế giới về file: $\mathrm{rot}'=F\cdot\mathrm{rot}$, $t'=F t$. Khi đó:
\begin{itemize}
\item vị trí $p'=(s\,\mathrm{rot}')p+t'$;
\item pháp tuyến $n'=\mathrm{rot}'\,n$ (xoay thuần, không co giãn);
\item scale (log): $\mathrm{scale}_i'=\mathrm{scale}_i+\ln s$, vì $\exp(\mathrm{scale}+\ln s)=s\exp(\mathrm{scale})$;
\item quaternion $q'=q_{\mathrm{rot}}\otimes q$;
\item SH: $\mathrm{sh}'=D(\mathrm{rot}')\,\mathrm{sh}$ (\S\ref{s:shrot}).
\end{itemize}

\subsection{\code{Layer.world_splats} (dùng cho vật lý, bounds)}
\code{positions = sp.positions @ (s*rot).T + t}; \code{rotations = quat_mul(q_rot, q)}; \code{scales *= s}.

%====================================================================
\section{Quaternion}\label{s:quat}
%====================================================================
\src{core/math3d.py}
Quaternion đơn vị $q=(w,x,y,z)$ biểu diễn phép quay góc $\theta$ quanh trục $n$: $w=\cos\frac\theta2$, $(x,y,z)=\sin\frac\theta2\,n$. Dùng quaternion vì không bị gimbal, nhân gọn và dễ chuẩn hóa.

\subsection{Quaternion sang ma trận (\code{quat_to_matrix}, giống hệt trong \code{splat_prep.vert})}
\[
R=\begin{bmatrix}
1-2(y^2+z^2) & 2(xy-wz) & 2(xz+wy)\\
2(xy+wz) & 1-2(x^2+z^2) & 2(yz-wx)\\
2(xz-wy) & 2(yz+wx) & 1-2(x^2+y^2)
\end{bmatrix}
\]
Trong GLSL, \code{mat3(...)} điền theo \textbf{cột}, nên các số liệt kê trong shader là cột của cùng ma trận này.

\subsection{Ma trận sang quaternion (\code{matrix_to_quat}, phương pháp Shepperd)}
Tránh chia cho số gần 0 bằng cách chọn nhánh theo phần tử lớn nhất.
\begin{itemize}
\item $\mathrm{trace}=R_{00}+R_{11}+R_{22}>0$: $s=2\sqrt{\mathrm{trace}+1}$, $w=s/4$, $x=(R_{21}-R_{12})/s$, $y=(R_{02}-R_{20})/s$, $z=(R_{10}-R_{01})/s$.
\item Ngược lại, chọn nhánh theo $R_{00}$, $R_{11}$ hoặc $R_{22}$ lớn nhất: $s=2\sqrt{1+R_{ii}-R_{jj}-R_{kk}}$, phần tử $i=s/4$, các phần tử còn lại $=(R_{ij}\pm R_{ji})/s$.
\end{itemize}

\subsection{Nhân quaternion (\code{quat_mul}, tích Hamilton)}
\begin{align*}
w&=a_wb_w-a_xb_x-a_yb_y-a_zb_z, & x&=a_wb_x+a_xb_w+a_yb_z-a_zb_y,\\
y&=a_wb_y-a_xb_z+a_yb_w+a_zb_x, & z&=a_wb_z+a_xb_y-a_yb_x+a_zb_w.
\end{align*}
$a\otimes b$ là quay $b$ trước rồi $a$ (giống $R_aR_b$).

\subsection{Lượng tử hóa quaternion lên GPU (\code{splat_renderer.pack}, \code{splat_prep.vert})}
Mỗi thành phần thuộc $[-1,1]$ được lưu 8 bit: $u=\mathrm{clip}(\mathrm{round}(q\cdot127)+128,1,255)$. Shader giải bằng $q=\mathrm{normalize}((u-128)/127)$, chuẩn hóa lại để bù sai số.

%====================================================================
\section{Camera và phép chiếu}\label{s:cam}
%====================================================================
\src{core/math3d.py, render/camera.py}

\subsection{Ma trận nhìn (\code{look_at(e, target, up)})}
\[
f=\frac{\mathrm{target}-e}{\|\cdot\|},\quad r=\frac{f\times \mathrm{up}}{\|\cdot\|},\quad u=r\times f,\quad
V=\begin{bmatrix} r^{\top} & -r\!\cdot\! e\\ u^{\top} & -u\!\cdot\! e\\ -f^{\top} & f\!\cdot\! e\\ 0\;0\;0 & 1\end{bmatrix}.
\]
Cột 3 là $-(Re)$: đưa mắt về gốc rồi xoay hệ trục camera về $(x,y,z)=(\text{phải},\text{lên},-\text{trước})$. Khối $3\times3$ trái trên gọi là $W$ (trực giao).

\subsection{Phép chiếu phối cảnh (\code{perspective(fov, a, n, f)})}
Với $k=1/\tan(\mathrm{fov}/2)$:
\[
P=\begin{bmatrix}k/a&0&0&0\\0&k&0&0\\0&0&\frac{f+n}{n-f}&\frac{2fn}{n-f}\\0&0&-1&0\end{bmatrix}
\]
\begin{itemize}
\item $\mathrm{clip}=Pv$, $c_w=-v_z$ (độ sâu dương). $\mathrm{NDC}=\mathrm{clip}_{xyz}/c_w\in[-1,1]$ (\textbf{chia phối cảnh}, vật xa nhỏ đi).
\item Hàng 3: $z_{ndc}=\dfrac{(f+n)v_z+2fn}{(n-f)(-v_z)}$ ánh xạ $v_z=-n\mapsto-1$ và $v_z=-f\mapsto+1$. Quan hệ này phi tuyến, độ phân giải tập trung gần camera.
\item Camera dùng \code{near = distance*0.001}, \code{far = near*1e6} (\code{camera.projection_matrix}).
\end{itemize}

\subsection{Phép chiếu trực giao (\code{orthographic(h, n, f)})}
$x'=x/h$, $y'=y/h$, $z'=-\dfrac{2z}{f-n}-\dfrac{f+n}{f-n}$. Hộp $[-h,h]^2\times[n,f]$ ánh xạ vào $[-1,1]^3$. Dùng cho \textbf{bản đồ bóng} mặt trời (tia song song, không chia phối cảnh). Hàm này bị định nghĩa trùng hai lần trong file, bản sau ghi đè bản trước, cùng công thức.

\subsection{Ví dụ số (\code{report/project3d_table.tex})}
$p_w=(0.5,1,0.5)\to v=Vp=(0.109,0.293,-2.536)\to \mathrm{clip}=Pv=(0.197,0.709,2.341,2.536)$, $c_w=2.536\to\mathrm{NDC}=(0.078,0.279,0.923)\to$ pixel $(215.5,108.1)$ theo \S\ref{s:viewport}.

\subsection{Camera quỹ đạo (\code{OrbitCamera})}
\begin{itemize}
\item \textbf{Vị trí mắt} (tọa độ cầu): $e=t+d\,(\cos\varphi\sin\theta,\;\sin\varphi,\;\cos\varphi\cos\theta)$ với $\theta$ là yaw, $\varphi$ là pitch, $d$ là distance (\code{position}).
\item Pitch kẹp $\pm1.55$ rad ($\approx\pm88.8^\circ$) để $f\times\mathrm{up}$ không suy biến (\code{PITCH_LIMIT}).
\item \code{orbit}: yaw $-=dx\cdot0.005$, pitch $+=dy\cdot0.005$ (rad/pixel).
\item \textbf{Pan} (\code{pan}): chiều cao nhìn thấy tại khoảng cách $d$ là $2d\tan(\mathrm{fov}/2)$, trải trên $H$ pixel, nên $\mathrm{units\_per\_pixel}=2d\tan(\mathrm{fov}/2)/H$. Khi đó $\mathrm{target}\mathrel{+}=(-dx\cdot\mathrm{right}+dy\cdot\mathrm{up})\cdot\mathrm{units\_per\_pixel}$, với right, up là hàng 0 và 1 của $V$.
\item \textbf{Look} (\code{look}): giữ mắt cố định, đổi yaw/pitch, rồi \code{target = eye - (position_moi - target)} để mắt không dịch.
\item \textbf{Bay} (\code{fly}): hướng ngang $=\mathrm{forward}\cdot(-\sin\theta,0,-\cos\theta)+\mathrm{right}\cdot(\cos\theta,0,-\sin\theta)$, chuẩn hóa (đi chéo không nhanh hơn), $y=\mathrm{up}$. Vận tốc mong muốn $v^*=\text{hướng}\cdot\text{tốc độ}\cdot(2.5\text{ nếu chạy})$. \textbf{Lọc mũ:} $v\leftarrow v+(v^*-v)\min(10\Delta t,1)$; $\mathrm{target}\mathrel{+}=v\Delta t$ (tích phân Euler). Dừng hẳn khi $|v|<10^{-3}\cdot\text{tốc độ}$.
\item \textbf{Frame} (\code{frame}): tâm là median các điểm; bán kính $R$ là phân vị 90\% khoảng cách tới tâm (bỏ nhiễu); $d=1.3R/\sin(\mathrm{fov}/2)$ (hình cầu bán kính $R$ vừa khung dọc: $R=d\sin(\mathrm{fov}/2)$); tốc độ bay $=0.5R$.
\item \textbf{Ngôi thứ ba:} mode 1: $\mathrm{cam}=\mathrm{eye}-\mathrm{forward}\cdot D$; mode 2: $\mathrm{cam}=\mathrm{eye}+\mathrm{forward}\cdot D$ nhìn ngược lại. $D$ co lại nếu tia mắt$\to$camera chạm voxel (\code{segment_blocked}, \S\ref{s:phys}).
\item $\mathrm{forward}()=(\mathrm{target}-\mathrm{position})/\mathrm{distance}$.
\end{itemize}

%====================================================================
\section{EWA Splatting: chiếu ellipsoid 3D thành ellipse 2D}\label{s:ewa}
%====================================================================
\src{render/shaders/splat_prep.vert}

\subsection{Hiệp phương sai 3D}
$\Sig = R S S^{\top} R^{\top} = M M^{\top}$, với $M=R\,\mathrm{diag}(s)$ (shader: \code{M = R * mat3(scale...)}, \code{sigma = M * transpose(M)}). $\Sig$ đối xứng bán xác định dương; trục ellipsoid là cột của $R$, bán trục là $s$.

\subsection{Jacobian của phép chiếu}
Điểm ảnh: $x_{px}=f_{px}v_x/(-v_z)$, $y_{px}=f_{px}v_y/(-v_z)$. Đạo hàm theo $(v_x,v_y,v_z)$ tại tâm:
\[
J=\begin{bmatrix}\frac{f}{-v_z}&0&\frac{f v_x}{v_z^2}\\0&\frac{f}{-v_z}&\frac{f v_y}{v_z^2}\\0&0&0\end{bmatrix},\qquad
f_{px}=\tfrac H2 P_{11}=\frac{H/2}{\tan(\mathrm{fov}/2)}
\]
(\code{splat_renderer}: \code{uFocal = height/2 * projection[1,1]}). Phép chia cho $c_w$ là phi tuyến, nên dùng \textbf{xấp xỉ tuyến tính bậc nhất} tại tâm (khai triển Taylor) để ellipsoid vẫn cho ra ellipse.

\subsection{Hiệp phương sai 2D}
\[
\Sig' = J\,W\,\Sig\,W^{\top}J^{\top}.
\]
Trong shader gộp thêm ma trận layer $sA$: \code{T = J * mat3(ModelView) * linear}, \code{cov = T Σ Tᵀ} ($\Sig$ ở tọa độ file; \code{linear = s·A} đã gồm co giãn nên không cần $s^2$ riêng). Lấy khối $2\times2$: $a=\mathrm{cov}_{00}$, $b=\mathrm{cov}_{01}$, $c=\mathrm{cov}_{11}$.

\textbf{Bộ lọc thông thấp:} $a\mathrel{+}=0.3$, $c\mathrel{+}=0.3$ (tương đương cộng $0.3I$), bảo đảm gaussian rộng ít nhất khoảng 1 pixel, chống răng cưa như EWA gốc và 3DGS.

\subsection{Trị riêng và vector riêng của ma trận 2x2}
\begin{align*}
\mathrm{mid}&=\tfrac{a+c}{2}, & \mathrm{radius}&=\sqrt{\left(\tfrac{a-c}{2}\right)^2+b^2},\\
\lambda_1&=\mathrm{mid}+\mathrm{radius}, & \lambda_2&=\max(\mathrm{mid}-\mathrm{radius},\,0.1)\ \text{(kẹp chống suy biến)},\\
\mathrm{axis}_1&=\mathrm{normalize}(b,\lambda_1-a), & \mathrm{axis}_2&=(-\mathrm{axis}_{1,y},\mathrm{axis}_{1,x}).
\end{align*}
Nếu $|b|<10^{-6}$ thì dùng trục $x$ hoặc $y$. Kiểm tra: $(A-\lambda_1 I)v=0\Rightarrow(a-\lambda_1)v_1+bv_2=0\Rightarrow v\propto(b,\lambda_1-a)$. Bán trục ellipse là $\sqrt{\lambda_1},\sqrt{\lambda_2}$ (độ lệch chuẩn theo hai hướng riêng).

\subsection{Kích thước hình vuông (billboard)}
\begin{itemize}
\item $\mathrm{extent}=\min\!\big(\sqrt{2\ln(255\alpha)},\,3\big)$, từ phương trình $\alpha e^{-r^2/2}=1/255\Rightarrow r^2=2\ln(255\alpha)$. Gaussian mờ cho hình vuông nhỏ; tối đa $3\sigma$.
\item $\mathrm{major}=\min(\mathrm{extent}\cdot\sqrt{\lambda_1},1024)\cdot\mathrm{axis}_1\cdot(2/\mathrm{viewport})$, \code{minor} tương tự với $\lambda_2$. Nhân $2/W$, $2/H$ để đổi pixel sang NDC (NDC rộng 2).
\item \textbf{Loại bỏ (cull):} $\alpha<1/255$; $c_w\le0$ (sau camera); $|\mathrm{ndc}_x|$ hoặc $|\mathrm{ndc}_y|>1.3$ (đệm 30\% vì ellipse có thể lấn vào màn hình).
\end{itemize}

%====================================================================
\section{Hiển thị gaussian và trộn alpha}\label{s:blend}
%====================================================================
\subsection{Pass 2: đỉnh (\code{splat.vert})}
Mỗi instance vẽ 4 góc $(\pm1,\pm1)$: \code{pos_ndc = center + cx*major + cy*minor}; \code{vCorner = (cx,cy)*extent} (đơn vị $\sigma$ trong hệ trục ellipse).

\subsection{Fragment (\code{splat.frag})}
\begin{itemize}
\item $r^2=|\mathrm{vCorner}|^2$; bỏ nếu $r^2>9$ (ngoài $3\sigma$).
\item $a=\alpha\exp(-r^2/2)$ (mật độ Gauss 2D chuẩn hóa về đỉnh bằng $\alpha$). Bỏ nếu $a<1/255$.
\item Xuất $(\mathrm{rgb}\cdot a,\,a)$, tức \textbf{premultiplied alpha}.
\end{itemize}

\subsection{Trộn ``over'' (\code{glBlendFunc(GL_ONE, GL_ONE_MINUS_SRC_ALPHA)})}
\[
C_i = c_i a_i + (1-a_i)\,C_{i-1}
\]
(với rgb đã nhân $a$: $C=\mathrm{src.rgb}+(1-\mathrm{src.a})\,\mathrm{dst}$). Phép này không giao hoán, nên phải vẽ \textbf{từ xa đến gần}. Depth test tắt (vật trong suốt không được che nhau).

\subsection{Biến đổi viewport (\code{brush._to_screen})}\label{s:viewport}
\[
x_{px}=x_0+\tfrac{n_x+1}{2}W,\qquad y_{px}=y_0+\tfrac{1-n_y}{2}H\ \text{(đảo $y$: màn hình có $y$ hướng xuống)}.
\]

\subsection{Tự chỉnh ngân sách (\code{SplatRenderer.adapt_budget})}
Đo thời gian GPU bằng query. $\mathrm{cost}=\mathrm{median}(\mathrm{ms}/\mathrm{count})$; $\mathrm{wanted}=\mathrm{target\_ms}/\mathrm{cost}$. Nếu lệch dưới 10\% thì giữ nguyên; ngược lại $\mathrm{step}=\mathrm{clip}(\mathrm{wanted}/\mathrm{current},0.75,1.33)^{1/2}$ (căn bậc hai nên đổi tối đa khoảng 15\%/lần, làm trơn để tránh dao động); $\mathrm{budget}=\mathrm{clip}(\mathrm{current}\cdot\mathrm{step},\,\mathrm{MIN\_BUDGET},\,\mathrm{count})$.

%====================================================================
\section{Lượng tử hóa và đóng gói GPU}
%====================================================================
\src{splat_renderer.pack, splat_prep.vert}
\begin{itemize}
\item \textbf{Half-float (IEEE-754 16 bit)}, \code{halfToFloat}: $e$ là bit 10--14, $m$ là bit 0--9, bit 15 là dấu. Nếu $e=0$ (subnormal) thì $v=m\cdot2^{-24}$; nếu $e\ne0$ thì $v=(m+1024)\cdot2^{e-25}=(1+m/1024)\cdot2^{e-15}$. Dùng cho màu, $\alpha$, scale để tiết kiệm bộ nhớ.
\item \textbf{SH 8 bit:} $\mathrm{scale\_sh}=\mathrm{percentile}_{99.9}(|\mathrm{coeff}|)/127$; lưu $u=\mathrm{clip}(\mathrm{round}(\mathrm{coeff}/\mathrm{scale\_sh})+128,1,255)$; giải $\mathrm{coeff}=(u-128)\cdot\mathrm{scale\_sh}$ (shader: \code{(texel*255 - 128) * uShScale}).
\item \textbf{Màu RGBA8} (\code{packColor}): $u=\mathrm{clamp}(c,0,1)\cdot255+0.5$ (làm tròn).
\item 1 gaussian $=8$ uint32 (2 texel RGBA32UI): xyz (float), màu rg/b+$\alpha$, scale xyz (half)+layer id, quaternion $4\times8$ bit.
\end{itemize}

%====================================================================
\section{Spherical Harmonics: màu theo hướng nhìn}
%====================================================================
\src{core/sh.py, splat_prep.vert: shColor}

\subsection{Công thức màu}
\[
c(d)=\mathrm{clip}\Big(0.5+C_0 f_{dc}+\sum_{k=1}^{15}Y_k(d)\,h_k,\;0,\;1\Big),
\]
với $d$ là vector đơn vị từ camera tới gaussian (3 kênh RGB độc lập).

\subsection{Hàm cơ sở thực (đúng quy ước dấu của 3DGS gốc)}
Với $d=(x,y,z)$, $xx=x^2$, \ldots:
\begin{itemize}
\item Bậc 1 (3 hàm): $-C_1y,\; C_1z,\; -C_1x$ với $C_1=\sqrt{3/4\pi}=0.4886025$.
\item Bậc 2 (5 hàm): $1.0925\,xy$, $-1.0925\,yz$, $0.3154\,(2zz-xx-yy)$, $-1.0925\,xz$, $0.5463\,(xx-yy)$. Các hằng số: $1.0925=\tfrac12\sqrt{15/\pi}$, $0.3154=\tfrac14\sqrt{5/\pi}$, $0.5463=\tfrac14\sqrt{15/\pi}$.
\item Bậc 3 (7 hàm): $-0.5900\,y(3xx-yy)$, $2.8906\,xyz$, $-0.4570\,y(4zz-xx-yy)$, $0.3732\,z(2zz-3xx-3yy)$, $-0.4570\,x(4zz-xx-yy)$, $1.4453\,z(xx-yy)$, $-0.5900\,x(xx-3yy)$.
\end{itemize}
Đây là các đa thức Legendre liên kết trên mặt cầu: $1+3+5+7=16=(L+1)^2$ với $L=3$. Số hệ số bậc cao $K\in\{0,3,8,15\}$ (\code{COEFFS_FOR_DEGREE}). Trong file, \code{f_rest_{c*K+k}} là hệ số $k$ của kênh $c$ (xếp theo kênh). Shader chọn bậc theo \code{uShDegree}.

\subsection{Xoay hệ số SH khi xoay gaussian (\code{sh.rotation_matrix}, \code{sh.rotate})}\label{s:shrot}
Khi gaussian quay $R$, hàm màu mới là $\mathrm{m\grave{a}u}'(d)=\mathrm{m\grave{a}u}(R^{\top}d)$. Mỗi bậc SH đóng kín dưới phép quay, nên tồn tại $D(R)$ ($K\times K$, khối chéo theo bậc) sao cho $\mathrm{sh}'=D\cdot\mathrm{sh}$. Thay vì dùng công thức Wigner-D, code \textbf{giải bình phương tối thiểu}: chọn 64 hướng $d_i$, cần $B(d_i)D=B(R^{\top}d_i)$, nên
\[
D=\mathrm{lstsq}\big(B(\mathrm{dirs}),\,B(\mathrm{dirs}\,R)\big)
\]
(hàng $d_i^{\top}R=(R^{\top}d_i)^{\top}$). Kết quả chính xác vì $B$ có hạng đầy đủ.
\begin{itemize}
\item \textbf{Hướng mẫu Fibonacci} (\code{_sample_dirs}): $z=1-2(i+\tfrac12)/n$, $r=\sqrt{1-z^2}$, $\varphi=i\,\pi(3-\sqrt5)$ (góc vàng $2.39996$ rad), cho điểm phân bố gần đều trên cầu.
\item Nhân khối bằng BLAS: $(n\cdot3,K)\,@\,D^{\top}$, xử lý từng \code{CHUNK} $2^{20}$.
\end{itemize}

\subsection{Trong shader (không xoay hệ số mỗi frame)}
$d_{\mathrm{file}}=\mathrm{normalize}(\mathrm{linear}^{\top}(p_w-e))$: đưa hướng nhìn \textbf{ngược} về tọa độ file (dùng $A^{\top}$ vì quay trực giao; scale dương bị chuẩn hóa mất). Cách này rẻ hơn xoay hệ số mỗi khung.

%====================================================================
\section{Ánh sáng Phong và bóng đổ}\label{s:phong}
%====================================================================
\src{splat_prep.vert: shade, shadowFactor; core/lighting.py}

\subsection{Pháp tuyến gaussian}
$n$ là trục ngắn nhất của ellipsoid: $\mathrm{axis}=\arg\min|s|$, $n=\mathrm{normalize}(\mathrm{linear}\cdot R[\mathrm{axis}])$ (cột $R$ tương ứng, đưa sang thế giới), lật nếu $n\cdot\mathrm{toCam}<0$. Với gaussian tròn (không có mặt phẳng): $\mathrm{roundness}=\mathrm{smoothstep}(0.35,0.9,s_{\min}/s_{\max})$, $n=\mathrm{normalize}(\mathrm{mix}(n,(0,1,0),\mathrm{roundness}))$. Hàm $\mathrm{smoothstep}(a,b,x)$: $t=\mathrm{clip}((x-a)/(b-a),0,1)$, kết quả $t^2(3-2t)$ (đa thức Hermite bậc 3).

\subsection{Phong}
\begin{align*}
w(n,l)&=\mathrm{clip}\!\left(\frac{n\cdot l+0.35}{1.35},0,1\right)\quad\text{(wrap diffuse: $n\cdot l=-0.35\mapsto0$, $1\mapsto1$)}\\
I_a&=c\,L_a\\
I_d&=c\sum_j L_j\,w(n,l_j)\,S_j\\
I_s&=k_s\sum_j L_j\max(r_j\cdot v,0)^m S_j,\quad r=\mathrm{reflect}(-l,n)=2(n\cdot l)n-l,\ \text{chỉ khi }n\cdot l>0\\
c_{\mathrm{out}}&=I_a+I_d+I_s\quad\text{(cuối cùng clamp 0..1)}
\end{align*}
$v=\mathrm{normalize}(\mathrm{camPos}-p)$; $m=$ \code{shininess} (mặc định 24); $k_s=$ \code{specular} (0.25).

\subsection{Đèn điểm: suy giảm}
$\mathrm{att}=\mathrm{clamp}(1-d^2/\mathrm{range}^2,0,1)^2$ (bình phương nên mượt về 0 tại $d=\mathrm{range}$). Tối đa 8 đèn, không đổ bóng. Màu $=$ màu $\times$ cường độ.

\subsection{Mặt trời (\code{Lighting.sun_direction})}
\begin{align*}
\mathrm{angle}&=\tfrac{\mathrm{hour}-6}{12}\pi\quad(\text{giờ }6\to0,\;12\to\tfrac\pi2,\;18\to\pi)\\
\mathrm{elevation}&=\sin(\mathrm{angle})\cdot65^\circ\quad(\text{âm ban đêm})\\
\mathrm{azimuth}&=\mathrm{azimuth}_0+(\mathrm{angle}-\tfrac\pi2)\\
\mathrm{dir}&=(\cos el\cos az,\;\sin el,\;\cos el\sin az)\quad(\text{cầu}\to\text{Descartes, đơn vị})
\end{align*}
Chu kỳ: \code{hour += dt*24/day_length (mod 24)}, \code{day_length = 36 s}, tức 1.5 s/giờ.
\begin{itemize}
\item $\mathrm{daylight}=\mathrm{smoothstep}(-0.08,0.3,\mathrm{dir}_y)$.
\item Màu mặt trời: $\mathrm{mix}((1,.5,.25),(1,.96,.88),\mathrm{smoothstep}(0,.35,s))\cdot I\cdot\mathrm{smoothstep}(-.02,.2,s)$ (cam lúc thấp, trắng lúc cao).
\item Ambient: $\mathrm{base}\cdot\mathrm{mix}((.12,.16,.32),1,\mathrm{daylight})$ (xanh tối ban đêm).
\item Bầu trời: $\mathrm{mix}(\text{đêm},\text{ngày},\mathrm{daylight})$ rồi trộn thêm cam với trọng số $0.55\exp(-(s/0.18)^2)$, một bump Gauss quanh đường chân trời (hoàng hôn).
\item $\mathrm{mix}(a,b,t)=a+(b-a)t$.
\end{itemize}

\subsection{Máy trạng thái ngày--đêm (\code{DayCycle})}
Mốc bắt đầu: DAWN 5h, DAY 7h, DUSK 17h, NIGHT 19h. \code{phase_for(h)}: $h\in[\mathrm{start},\mathrm{end})$, hoặc nếu $\mathrm{start}>\mathrm{end}$ (qua nửa đêm) thì $h\ge\mathrm{start}\lor h<\mathrm{end}$. Sự kiện \code{on_enter} bật/tắt đèn \code{auto_night}.

\subsection{Bản đồ bóng (\code{Lighting.shadow_setup}, \code{shadow.vert}, \code{shadowFactor})}
\begin{itemize}
\item $r=1.1\cdot\mathrm{radius}$ cảnh; mắt đèn $\mathrm{eye}=\mathrm{center}+d\cdot2r$; $\mathrm{up}=(0,1,0)$ (hoặc $(0,0,1)$ nếu $|d_y|\ge0.99$ để tránh $f\parallel\mathrm{up}$).
\item $VP=\mathrm{Ortho}(r,0.01r,4r)\cdot\mathrm{LookAt}(\mathrm{eye},\mathrm{center})$; độ phân giải $2048^2$.
\item $\mathrm{texel}=2r/2048$; $\mathrm{bias}=3\,\mathrm{texel}/(3.99r)$ (độ lệch độ sâu quy về NDC $[0,1]$; miền $z$ dài $3.99r$).
\item Điểm tra: $lp=VP\cdot(p+n\cdot2\,\mathrm{texel})$ (đẩy dọc pháp tuyến chống tự bóng); $s=lp_{xyz}/lp_w\cdot0.5+0.5$ (NDC sang $[0,1]$ để lấy mẫu texture); nằm ngoài $[0,1]$ thì coi là sáng.
\item \textbf{PCF $3\times3$:} $S=\frac19\sum_{i,j=-1}^{1}\big[s_z-\mathrm{bias}\le D(s_{xy}+(i,j)/2048)\big]$.
\item Chỉ gaussian có $\alpha\ge0.3$ và chưa xóa mới được vẽ vào depth map (dạng điểm).
\end{itemize}

%====================================================================
\section{Texture}
%====================================================================
\begin{itemize}
\item \textbf{Skin $64\times64$} kiểu chữ thập Minecraft, \code{avatar_renderer._uv_region(face, ox, oy, w, h, d)} ($w,h,d$ là rộng, cao, sâu của hộp). Mặt trước $(ox+d,\,oy+d)$; sau $(ox+2d+w,\,oy+d)$; phải $(ox,\,oy+d,\,d\times h)$; trái $(ox+d+w,\,oy+d)$; trên $(ox+d,\,oy)$; dưới $(ox+d+w,\,oy)$. $\mathrm{UV}=\mathrm{pixel}/64$. Bốn đỉnh mỗi mặt: $(r_x,r_y)$, $(r_x+r_w,r_y)$, $(r_x+r_w,r_y+r_h)$, $(r_x,r_y+r_h)$.
\item Hộp: nửa kích thước $h_x=w/2+\mathrm{inflate}$, \ldots (lớp áo ``inflate'' phồng ra để tránh z-fighting). Đỉnh $=$ tâm $+(\pm1)\cdot$nửa kích thước; pháp tuyến là 6 hướng trục.
\item \textbf{Avatar shader:} $\mathrm{shade}=\mathrm{bright}\cdot(0.62+0.38\max(n\cdot l,0))$ (ambient 0.62 + Lambert 0.38); $n_{\mathrm{world}}=\mathrm{normalize}(\mathrm{mat3}(M)\,n)$; bỏ texel $\alpha<0.5$ (lớp áo trong suốt).
\item \textbf{Giải mã PNG tự viết} (\code{load_png}): zlib + bộ lọc dòng, với $a$ là trái, $b$ là trên, $c$ là trên-trái và mọi phép cộng đều mod 256: 1 Sub: $x+a$; 2 Up: $x+b$; 3 Average: $x+\lfloor(a+b)/2\rfloor$; 4 \textbf{Paeth}: $p=a+b-c$, chọn trong $\{a,b,c\}$ gần $p$ nhất ($p_a=|p-a|$, $p_b$, $p_c$; ưu tiên $a\le b\le c$).
\item Texel buffer: dữ liệu gaussian đọc bằng \code{texelFetch} (số nguyên, không nội suy).
\end{itemize}

%====================================================================
\section{Sắp xếp theo độ sâu}\label{s:sort}
%====================================================================
\src{render/sorter.py: DepthSorter.sort / _depth_rows}
\begin{itemize}
\item Độ sâu camera của tâm: $z_{cam}=(MV\cdot M_{layer})[2,:]\cdot(p,1)$. Camera nhìn $-z$ nên $z_{cam}<0$ là phía trước, càng âm càng xa. Giữ $\mathrm{depth}<0$.
\item Khóa 16 bit: $\mathrm{key}=\lfloor(\mathrm{depth}-lo)\cdot65535/(hi-lo)\rfloor$ (lượng tử đều trên $[lo,hi]$); \code{argsort(uint16, stable)} là radix sort $O(N)$. Depth (âm) tăng dần tương ứng xa $\to$ gần. Lượng tử 16 bit đủ vì chỉ cần thứ tự gần đúng.
\item Chạy ở luồng nền; kết quả cũ bị bỏ nếu \code{version} đã đổi.
\end{itemize}

%====================================================================
\section{LOD, Morton, Frustum}\label{s:lod}
%====================================================================
\src{core/lod.py}

\subsection{Mã Morton 3D, đường cong Z (\code{morton_codes})}
Hộp bao theo phân vị 0.1\%--99.9\% (bỏ nhiễu): $q=\mathrm{clip}((p-lo)\cdot1023/(hi-lo),0,1023)$ (10 bit/trục). Xen kẽ bit: \code{spread(v)} đặt bit $i$ của $v$ vào vị trí $3i$ (bảng tra); $\mathrm{code}=\mathrm{spread}(q_x)\mid\mathrm{spread}(q_y)\ll1\mid\mathrm{spread}(q_z)\ll2$ (30 bit). Điểm gần nhau trong không gian có mã gần nhau, tốt cho cache.

\subsection{Radix argsort (\code{radix_argsort})}
Hai lượt ổn định 16 bit (thấp rồi cao) cho thứ tự đúng theo khóa 32 bit (LSD radix).

\subsection{Khối và độ quan trọng (\code{_build_group})}
Cắt theo Morton thành khối 1024. Trong khối xếp theo $\mathrm{importance}=\alpha\,(s_0s_1+s_1s_2+s_0s_2)$ giảm dần (tỉ lệ với diện tích bề mặt ellipsoid $\times$ độ đục). Mẹo: khóa $=65535-(16\text{ bit cao của float32 dương})$, vì float dương sắp theo thứ tự bit nên không cần log.

\subsection{Cầu bao mỗi khối}
$\mathrm{radius}=\tfrac12\|hi-lo\|+3\max(s)$ (nửa đường chéo hộp tâm $+3\sigma$ lớn nhất). Khi có ma trận layer: $\mathrm{center}'=A\,\mathrm{center}+t$, $\mathrm{radius}'=\mathrm{radius}\cdot\|\text{cột 0 của }M\|=\mathrm{radius}\cdot s$.

\subsection{Frustum culling (\code{_frustum_planes}, \code{select})}
Với $VP=PV$, 6 mặt phẳng (Gribb--Hartmann): hàng~3 $\pm$ hàng~0/1/2 cho $(a,b,c,d)$, chuẩn hóa bằng $\|(a,b,c)\|$. Khoảng cách có dấu $\mathrm{dist}=a\,c_x+b\,c_y+c\,c_z+d$; khối nhìn thấy nếu $\mathrm{dist}>-\mathrm{radius}$ với 5 mặt (bỏ mặt xa).

\subsection{Phân bổ ngân sách (\code{select})}
Khối $c$ lấy $\mathrm{take}_c=\min(n_c,\lceil\lambda\,n_c/d_c\rceil)$ với $d_c=\max(\mathrm{depth}-\mathrm{radius},\epsilon)$. Tìm $\lambda$ bằng \textbf{chia đôi 40 vòng} sao cho $\sum\mathrm{take}=\mathrm{budget}$. Khối gần ($d$ nhỏ) lấy nhiều gaussian, khối xa lấy ít; ưu tiên top importance vì đã được xếp trước. Độ sâu khối: $\mathrm{depth}=-(\mathrm{center}\cdot mv_2+mv_{23})$.

%====================================================================
\section{Cọ chọn}\label{s:brush}
%====================================================================
\src{editor/brush.py}

\subsection{Khoảng cách điểm--đoạn, capsule (\code{_segment_dist2})}
$\tau=\mathrm{clip}\!\big((p-a)\cdot(b-a)/\|b-a\|^2,0,1\big)$, $d^2=\|p-(a+\tau(b-a))\|^2$; chọn nếu $d^2\le r^2$. Nếu $a=b$ thì $\tau=0$. Capsule là tập điểm cách đoạn không quá $r$ (hình viên thuốc, lồi).

\subsection{Chiếu tâm ra pixel}
$\mathrm{clip}=M_{[:3,:3]}p+M_{[:3,3]}$ với $M=VP\cdot M_{layer}$; chỉ xét $w>10^{-6}$ (trước camera); $x=\mathrm{left}+(c_x/w+1)/2\cdot W$, $y=\mathrm{top}+(1-c_y/w)/2\cdot H$.

\subsection{Tăng tốc theo khối}
8 đỉnh hộp bao khối (\code{_CORNERS}: bit $i$ của chỉ số chọn lo/hi) được chiếu thành hình chữ nhật bao (hộp lồi và phép chiếu bảo toàn tính lồi khi $w>0$). Khối chạm nét nếu hình chữ nhật giao hộp $[\min(a,b)-r,\max(a,b)+r]$. Khối \textbf{nằm trọn trong nét} nếu cả 4 góc hình chữ nhật có $d^2\le r^2$ (capsule lồi nên cả hình chữ nhật nằm trong), khi đó lấy cả khối mà không kiểm từng gaussian. Khối cắt mặt phẳng camera (dấu $w$ lẫn lộn) luôn được kiểm. Cọ \textbf{xuyên sâu} (không test độ sâu) như SuperSplat.

%====================================================================
\section{Vật lý và va chạm}\label{s:phys}
%====================================================================
\src{core/physics.py}

\subsection{Voxel hóa (\code{CollisionWorld.build})}
$\mathrm{cell}=\lfloor(p-\mathrm{origin})/v\rfloor$, $\mathrm{origin}=\min-v$; $v=\mathrm{height}/8$; kẹp $[0,2^{21}-1]$. \textbf{Khóa 63 bit:} $\mathrm{key}=(c_x\ll42)\mid(c_z\ll21)\mid c_y$ (mỗi trục 21 bit). $c_y$ ở bit thấp nên các ô cùng cột $(c_x,c_z)$ nằm liền nhau trong mảng đã sắp. Ô đặc nếu có ít nhất \code{min_count} gaussian (3/2/1 tùy mức lấy mẫu). \code{floor_y} là phân vị 1\% của $y$ (sàn an toàn).

\subsection{Truy vấn cột (\code{_highest})}
$\mathrm{query}=\mathrm{key}(i_x,i_z,i_y)$; $\mathrm{idx}=\mathrm{searchsorted}(\mathrm{keys},\mathrm{query},\text{`right'})-1$, tức khóa lớn nhất $\le\mathrm{query}$; hợp lệ nếu $\mathrm{found}\gg21=\mathrm{query}\gg21$ (cùng cột). Trả về $i_y$ cao nhất $\le y$. Độ cao mặt trên ô: $\mathrm{origin}_y+(i_y+1)v$. Độ phức tạp $O(\log N)$ mỗi truy vấn, đã vector hóa.

\subsection{Tia kiểm tra tầm nhìn (\code{segment_blocked})}
Lấy mẫu dọc đoạn với bước $v/2$: $n=\lfloor L/(v/2)\rfloor+1$, $t=\mathrm{linspace}(0,1,n)$; điểm $=\mathrm{start}+(\mathrm{end}-\mathrm{start})t$; tra khóa; trúng ô đặc thì bị chặn.

\subsection{Người chơi}
Kích thước tỉ lệ với chiều cao $H$ (2.74): mắt $0.92H$, bán kính $0.14H$, bước leo $0.35H$, nhảy $0.75H$, $g=12H$, tốc độ đi $2.5H$ ($\times1.7$ chạy, $\times0.4$ rình).
\begin{itemize}
\item \textbf{Va chạm cột trụ:} 9 điểm vòng \code{_RING} $=$ tâm + 8 điểm $(\cos\frac{k\pi}{4},\sin\frac{k\pi}{4})\cdot\mathrm{bán\ kính}$.
\item \textbf{Trượt tường:} thử $(x+dx,z+dz)$; nếu bị chặn thử chỉ $x$; rồi chỉ $z$.
\item \textbf{Nhảy:} $v_y=\sqrt{2g\,h_{\mathrm{nhảy}}}$ (từ $v^2=2gh$).
\item \textbf{Euler:} $v_y\leftarrow v_y-g\Delta t$, $y\leftarrow y+v_y\Delta t$ ($\Delta t$ kẹp $\le0.05$); chạm đất khi $v_y\le0\ \land\ y_{\mathrm{mới}}\le\mathrm{ground}$.
\item \textbf{Làm mượt mắt:} $\mathrm{eye}_y\mathrel{+}=(\mathrm{target}-\mathrm{eye}_y)\min(18\Delta t,1)$ (lọc mũ, che bậc thang).
\end{itemize}

%====================================================================
\section{Hoạt họa nhân vật}\label{s:anim}
%====================================================================
\src{core/avatar_pose.py}
\begin{itemize}
\item Bộ lọc tiến về mục tiêu: $\mathrm{approach}(v,t,\mathrm{rate},dt)=v+(t-v)\min(1,dt\cdot\mathrm{rate})$ (cùng dạng quán tính camera bay).
\item \textbf{Nhịp đi:} $\mathrm{walked}=\mathrm{distance}/\mathrm{height}$; $\mathrm{phase}\mathrel{+}=\mathrm{walked}\cdot\mathrm{STRIDE}\cdot1.75$ (pha tỉ lệ \textbf{quãng đường}, không phải thời gian, nên chân không trượt); $\mathrm{swing}=\sin(\mathrm{phase})\cdot60^\circ\cdot\mathrm{amount}$; chân trái $+\mathrm{swing}$, chân phải $-\mathrm{swing}$, tay ngược pha (đối xứng chéo).
\item Biên độ mục tiêu $=\min(1,\mathrm{speed}/(4.3/1.75))$.
\item \textbf{Thở:} $\mathrm{idle}=2.2\sin(1.13t)+0.8\sin(1.91t)$ (hai tần số không công ước nên không lặp cứng).
\item \textbf{Cử chỉ:} zombie $-85^\circ+3\sin(2.4t)$; cheer $\pm(158+6\sin7t)$; wave $142+20\sin9t$; pha trộn tuyến tính \code{_mix_pair}.
\item Ngồi: chân $-90^\circ$, hạ thân 10 đơn vị skin; $\mathrm{stance}=1-0.48\,\mathrm{sit}-0.12\,\mathrm{crouch}$.
\end{itemize}

\subsection{Cây phân cấp (\code{limb_matrices})}
\begin{align*}
\mathrm{root}&=T(\mathrm{feet})\,R_y(\mathrm{facing}+180^\circ)\,S(\mathrm{scale})\,T(0,-\mathrm{drop},0)\\
\mathrm{upper}&=\mathrm{root}\cdot T(0,12,0)\,R_x(\mathrm{lean})\,T(0,-12,0)\quad(\text{xoay quanh hông})\\
\mathrm{leg}&=\mathrm{root}\cdot T(0,12,0)\,R_x(a)\,T(0,-12,0)\\
\mathrm{arm}&=\mathrm{upper}\cdot T(p_x,24,0)\,R_z(a_1)\,R_x(a_0)\,T(-p_x,-24,0)\quad(\text{quanh vai, }p_x=\pm6)\\
\mathrm{head}&=\mathrm{upper}\cdot T(0,24,0)\,R_y(\mathrm{yaw})\,R_x(\mathrm{pitch})\,T(0,-24,0)\quad(\text{quanh cổ})
\end{align*}
Mẫu chung: $M_{\mathrm{con}}=M_{\mathrm{cha}}\cdot T(p)\,R\,T(-p)$ (xoay quanh điểm $p$: dời $p$ về gốc, xoay, dời lại). $\mathrm{scale}=\mathrm{chiều\ cao}/32$. Ma trận quay trục \code{_rot(axis, deg)} điền $c,-s,s,c$ vào cặp chỉ số $(1,2)$ cho $X$, $(2,0)$ cho $Y$, $(0,1)$ cho $Z$.

%====================================================================
\section{Hình học các thành phần phụ}
%====================================================================
\subsection{Lưới sàn (\code{grid.frag})}
Tia từ camera: $\mathrm{dir}=\mathrm{CamRot}\cdot(\mathrm{ndc}_x/P_{00},\,\mathrm{ndc}_y/P_{11},\,-1)$ (ngược phép chiếu). \textbf{Giao tia--mặt phẳng $y=h$:} $t=(h-e_y)/\mathrm{dir}_y$, $p=e+t\cdot\mathrm{dir}$ (bỏ nếu $|\mathrm{dir}_y|<10^{-6}$ hoặc $t\le0$).

Đường lưới: $c=p/\mathrm{cell}$, $w=\mathrm{fwidth}(c)$ (đạo hàm màn hình, tức độ rộng 1 pixel), $g=|\mathrm{fract}(c-\tfrac12)-\tfrac12|/w$, $\mathrm{line}=1-\min(\min(g_x,g_y),1)$ (khử răng cưa theo đạo hàm). Mức chi tiết liên tục: $\mathrm{level}=\log_{10}(\mathrm{size})$, $\mathrm{cell}=10^{\lfloor\mathrm{level}\rfloor-1}$, $\mathrm{blend}=\mathrm{level}-\lfloor\mathrm{level}\rfloor$; lưới lớn $=10\,\mathrm{cell}$. Mờ theo khoảng cách $1-\mathrm{smoothstep}(\mathrm{fade},3\,\mathrm{fade},\mathrm{dist})$ và theo góc nhìn $\mathrm{smoothstep}(0,.15,|\mathrm{dir}_y|)$.

\subsection{Tam giác phủ toàn màn hình (\code{grid.vert})}
$p=((\mathrm{id}\ll1)\,\&\,2,\;\mathrm{id}\,\&\,2)\cdot2-1$, nên $\mathrm{id}=0,1,2$ cho $(-1,-1)$, $(3,-1)$, $(-1,3)$.

\subsection{Đom đóm (\code{firefly.vert})}
Góc vàng $i\cdot2.399963$ ($=\pi(3-\sqrt5)$) cho pha. Offset $=(\sin(\omega t+\varphi),\;0.6\sin(0.7t+1.3\varphi)+0.2,\;\cos(\omega't+0.7\varphi))\cdot R\cdot(0.35+0.65\,\mathrm{fract}(0.618\,i))$, trong đó $0.618$ là nghịch đảo tỉ lệ vàng (dãy Weyl phân bố đều bán kính). Nhấp nháy $0.5+0.5\sin(2t+3\varphi)$. Điểm mờ: $a=\mathrm{clamp}(1-d,0,1)^2\cdot\mathrm{twinkle}$.

\subsection{Quả cầu đèn (\code{light_gizmo.*})}
Billboard: \code{view.xy += corner*R*halo}. Trên hình tròn $r\le1$: $z=\sqrt{1-r^2}$ (cao độ mặt cầu), cho độ sáng $0.7+0.5z$. Bên ngoài là quầng $h=\big(1-\frac{r-1}{\mathrm{halo}-1}\big)^2\cdot0.65$. Bán kính hiển thị $\max(0.12\,\mathrm{unit},\,0.012\,\mathrm{khoảng\ cách})$ (đủ lớn trên màn hình).

Vị trí camera từ view: $e=-R^{\top}t$ (\code{splat_renderer}, \code{light_renderer}, \code{grid_renderer}).

%====================================================================
\section{Undo}
%====================================================================
\src{core/history.py}
Mẫu Command, tối đa 20 bước. Xóa mềm bằng bit \code{deleted}; hoàn tác là lật lại các chỉ số ($O(k)$). Trạng thái GPU 1 byte/gaussian: bit 0 = chọn, bit 1 = xóa (\code{STATE_*}). Màu chọn: $\mathrm{mix}(\mathrm{color},\mathrm{selectColor},0.7\cdot\mathrm{selected})$.

%====================================================================
\section{Bảng tra nhanh: công thức và file}
%====================================================================
\begin{center}\small
\begin{tabular}{@{}p{0.42\textwidth}p{0.5\textwidth}@{}}\toprule
Chủ đề & File:hàm\\\midrule
sigmoid, exp, $C_0$, chuẩn hóa quaternion & \code{io/codec.py: _decode_block, encode}\\
Euler, Rodrigues, quaternion & \code{core/math3d.py}\\
$M=[sA\,|\,t]$, xoay quanh pivot & \code{core/layer.py: world_matrix}; \code{core/transform.py: rotate_about}\\
view, projection, ortho & \code{core/math3d.py: look_at, perspective, orthographic}\\
orbit, pan, fly, frame & \code{render/camera.py}\\
EWA $\Sig'=JW\Sig W^{\top}J^{\top}$, eigen, extent & \code{render/shaders/splat_prep.vert: main}\\
Gauss và trộn over & \code{render/shaders/splat.frag}; \code{splat_renderer.py: _draw_splats}\\
SH cơ sở và xoay & \code{core/sh.py}; \code{splat_prep.vert: shColor}\\
Phong, suy giảm, wrap & \code{splat_prep.vert: shade}\\
Mặt trời, ngày--đêm, bóng & \code{core/lighting.py}\\
PCF & \code{splat_prep.vert: shadowFactor}\\
Morton, radix, frustum, budget & \code{core/lod.py}\\
Sắp xếp độ sâu 16 bit & \code{render/sorter.py}\\
Cọ capsule & \code{editor/brush.py}\\
Voxel, khóa 63 bit, trọng lực & \code{core/physics.py}\\
Hoạt họa và ma trận khớp & \code{core/avatar_pose.py}\\
Grid ray--plane, fwidth & \code{render/shaders/grid.frag}; \code{grid_renderer.py}\\
Skin UV, Paeth PNG & \code{render/avatar_renderer.py}\\
Half-float, lượng tử hóa & \code{splat_renderer.py: pack}; shaders\\\bottomrule
\end{tabular}
\end{center}

